\chapter{Ethics of Doxonomics}
\label{ch:ethics}

\epigraph{The question is not whether we will be extremists, but what kind of extremists we will be.}{Martin Luther King Jr.}

\noindent
A science that studies how beliefs are manufactured must confront an uncomfortable question: does studying belief engineering risk becoming belief engineering?
This chapter addresses the ethical tensions inherent in doxonomics.

\section{Description vs.\ Intervention}
\label{sec:ch34-description}

Doxonomics is primarily descriptive: it maps how beliefs are produced, not how they should be produced.
But description is never fully neutral.
Revealing that a belief is heavily subsidized changes its status.
A doxonomic audit is itself a doxonomic intervention.

\begin{principle}[Transparency Principle]
\label{prin:transparency}
\index{transparency principle}
Doxonomic research should be transparent about its own funding, institutional position, and potential biases.
The field that studies how credibility is manufactured must maintain its own credibility through radical openness.
\end{principle}

\section{The Danger of Technocratic Belief Management}
\label{sec:ch34-technocracy}

\begin{definition}[Doxonomic Technocracy]
\label{def:doxonomic-technocracy}
\index{doxonomic technocracy}
A \emph{doxonomic technocracy} is a regime in which experts use doxonomic tools to manage public belief ``for the public good.''
This is the dark twin of epistemic democracy: belief management by experts who decide what people should think.
\end{definition}

Doxonomics must resist this temptation.
The purpose of the field is to make belief production \emph{visible and accountable}, not to replace one set of belief engineers with another.

\section{Research Ethics}
\label{sec:ch34-research}

Doxonomic research raises specific ethical issues:
\begin{itemize}[nosep]
  \item experiments that manipulate beliefs may have lasting effects,
  \item revealing funding sources may endanger whistleblowers,
  \item doxonomic analysis of specific institutions may be perceived as political attack,
  \item the tools of doxonomics can be used by authoritarian regimes to refine propaganda.
\end{itemize}

\section{The Paradox}
\label{sec:ch34-paradox}

\begin{remark}[The Doxonomic Paradox]
\label{rem:paradox}
\index{doxonomic paradox}
Doxonomics studies how power shapes belief, but doxonomics itself is shaped by the power structures within which it operates.
Its funding, institutional home, and audience determine what questions it can ask and what answers it can publish.
No doxonomic analysis is exempt from doxonomic analysis.
\end{remark}

The field must therefore be reflexive: it must apply its own tools to itself, disclosing its funding, biases, and institutional constraints.
This is not a weakness but a strength---a science that audits itself is harder to co-opt.

\section{Notes and References}
\label{sec:ch34-notes}

On reflexivity in social science, see \textcite{bourdieu1984}.
The ethics of social research are discussed in \textcite{salganik2018}.
On the dangers of technocratic governance, see \textcite{scott1998} and \textcite{habermas1962}.
The dual-use problem in research connects to \textcite{oreskes2010}.

\section{Exercises}

\begin{exercise}
Apply doxonomic analysis to this textbook.
Who funded its production?
What institutional context shapes its narrative?
What beliefs does it promote?
\end{exercise}

\begin{exercise}
Design an ethics protocol for doxonomic research that balances transparency with the protection of sources and subjects.
\end{exercise}

\begin{exercise}
Argue for or against the following proposition: ``Doxonomic tools should be classified as dual-use technology, subject to restrictions on their application by state propaganda agencies.''
\end{exercise}
